\textbf{Lisboa} é um ponto importante para compreender as relações históricas entre Portugal e o Magrebe. A cidade foi marcada pela presença islâmica, pelas conexões comerciais com o Norte da África e pela expansão portuguesa pelo Atlântico. Antes do Marrocos, vale explorar suas camadas medieval, judaica, islâmica e marítima --- e perceber como essas histórias se prolongam do outro lado do Estreito de Gibraltar.

\textbf{Marraquexe} foi fundada pelos almorávidas no século \textsc{xi} e tornou-se uma das grandes capitais do Ocidente islâmico. Sua medina reúne palácios, mesquitas, madrasas, jardins e souks, com forte influência das tradições do Magrebe e de Al-Andalus. A cidade também preserva importantes monumentos saaditas e alauítas, além do Mellah e da praça Jemaa el-Fna. 

\textbf{Essaouira} é a antiga Mogador, uma cidade portuária atlântica planejada no século \textsc{xviii} pelo sultão Sidi Mohammed ben Abdallah para ampliar o comércio do Marrocos com a Europa. Sua medina combina arquitetura marroquina e fortificações de inspiração europeia. Porto cosmopolita, reuniu comunidades muçulmanas, judaicas, cristãs, amazigh e europeias.

\textbf{Fès} é um dos principais centros históricos, religiosos e intelectuais do Marrocos. Fundada no século \textsc{viii}, sua medina preserva mesquitas, madrasas, mercados e oficinas que refletem séculos de cultura islâmica e influência andaluza. A cidade também abriga uma importante memória judaica, especialmente no Mellah, e mantém uma forte tradição artesanal. 


